\section{Regression Machine Learning Model}
The advance of machine learning has led the wide applications in IoT devices.

KNNR: works by comparing a node with the surrounding nodes to predict a value. In a sense, KNNR is considered a lazy model, in which it remembers all training data, thus result in a heavy model

Gaussian Process Regressor, similar to SVM. However, in \cite{}, the authors analyse GPR and conlcude that the memory requirement of GPR is $O(n^2)$ while the computation cost is $O(n^3)$, which is quite computational and memory expensive for constrained devices

Deep learning is recently a new adavnces and has achieve in terms of matching the complex patterns and data. However, due to the constrained environment of the working MCU which has several KB of RAM and ROM. I rule out this.

Random forest and Decision Tree operates on similar algorithm. In tree base algorithm, each value is split into left and right, following the path to reach the end. In Random Forest, multiple trees are used to prevent overfitting and provide better performance. However, in my experiments, the performance of random forest and decision tree are quite similar to each other. Thus I remove random forest from list of considered models as they are much heavier than decision tree.
